\chapter{学工社团}
\section{学工简介}
大学生活的一个重要组成部分就是学生工作。

\vspace{10pt}

区别于中学，在大学中做学生工作的同学们有更多可以调动的资源，当然这也就意味着更大的责任。如果说“努力学习”是在学术方向上的努力，那么学生工作就是在社会方向上的探索。它教会你怎么服务更广大的群体，怎样和人打交道，怎样做出让大家都满意的工作，所以学生工作是服务他人的方式，更是锻炼自己的机会。

\vspace{10pt}

学生工作范畴极大，大体工作领域可以分为：班级、院学生会、院团委、校学生会、校团委、学校各级工作单位。

\vspace{10pt}

班级干部工作，主要负责上传下达，跟中学类似，不做赘述。

\vspace{10pt}

接下来简要介绍院级单位和校级单位工作区别。在院级单位，无论是学生会还是团委，我们的工作绝大多数都是围绕学院的师生展开的。优势是我们可以和同院师生更加亲近，能收获更好的关怀；劣势是很少走出学院的“一亩三分地”，活动规模相对较小。在校学生会和校团委的工作则面向全校师生展开，可以结识各个学院的同学老师，参与组织学校的大规模活动，但是可能和学院的同学相对比较生疏。

\vspace{10pt}

学生会和团委（院、校同理）的区别，是很多同学都会有的困惑。简要介绍如下：团委举办的活动相对更加正式，强调思政引领；工作模式和风格相对更加规范，组织严明。学生会是党政机关和广大学生联络的纽带，其工作内容和团委有一定程度的交叉，组织架构和团委也较为类似，但是举办的活动相对轻松，工作模式和风格相对自由。

\vspace{10pt}

是院团委、院会，还是校团委、校会，如何选择需要同学们自行衡量。

\vspace{10pt}

最后我们来说学校各级工作单位的学生助理。很多单位都招收学生助理，例如学生工作办公室助理，招生办公室助理……这些工作的工作时间固定，需要跟老师协调；工作内容是协助老师完成日常工作，对于了解该单位的工作以及相关领域的动向有一定帮助；另外值得一提的是，这些工作有固定的薪资酬劳。工作的招募一般在该单位的公众号上发布，有意的同学可以密切关注公众号。

\vspace{10pt}

同学们还可以通过辅导员老师，或者通过学院官网coe.pku.edu.cn——教职员工——行政人员，联系各位老师解决学生工作方面的问题。

\section{学工组织}

有关校学生会、校团委，上面已经有过一些介绍。详情请关注公众号“北京大学学生会”和“北大团委”，在这里不做赘述。下面详细介绍工学院学生会和团委的情况。

\vspace{10pt}

\subsection{工学院团委}

在工学院团委这个温馨的大家庭里，你可以收获：

\begin{itemize}
    \item 深度参与学院大型活动组织策划的体验。
    
    \item 更早获悉学院一手信息的机会。
    
    \item 较为专业的办公技能锻炼以及组织能力磨炼的机会。
    
    \item 丰富的阅历和视野，同时通过社会实践/志愿服务团体（例如组织参加支教等志愿服务活动）增长社会技能。
    
    \item 有优秀的学长学姐耐心引路，与志同道合的朋友共同进步的机会。
\end{itemize}

\vspace{10pt}

学院团委主要包含办公室、组织部、宣传部、实践部、社团运营中心、文体部、体育发展中心、青年志愿者协会、周培源宣讲团等部门，招新时间待定，招新对象为工学院全体同学。详细介绍请参考公众号“工映青春”上的招新推送。去年的推送为2024.9.11《招新｜工学院团委部门招新啦！》，敬请期待今年的招新推送。

\vspace{10pt}

补充：学生工作不是所有同学一定要参加的。事实上，北大任何活动几乎都没有强制性。但参加学生工作，不管是在功利角度还是非功利角度，都会有莫大的收获。请同学们根据自身实际情况，尽力而为，量力而行。

\vspace{10pt}


\subsection{工学院学生会}

学生会是联系学校、联系广大同学的桥梁和纽带。工学院学生会的宗旨是服务同学、虽小却精、氛围融洽、活动自主。

\vspace{10pt}

在这里你能体验到：

\begin{itemize}
    \item 有趣而温暖的文化。学生会以友爱为先，工作氛围极好，工作压力较小，没有恶性竞争，能获得很强的归属感。
    
    \item 丰富而自主的活动。不仅可以体验到各种活动，更可以提供你的点子，为学生会和学院同学创办一些新的活动方案。相对而言，团委的工作更偏向于正规和纪律性，请大家选择自己合适的工作风格。
    
    \item 满满当当的成就感。活动能锻炼同学们的能力和眼界，在学生会自主性质的加持下，更能让同学们体会到成功举办活动带来的满满当当的成就感和自我实现感。例如，本手册就是完全由工院学生会学术实践部自主发起和编写的，在撰写的过程中编者也感受到了自我价值的实现！
    
    \item 优秀而贴心的学长学姐。学生会的学长学姐不仅本身优秀，更是极其愿意帮助低年级的同学们成长向上，在这里你会收获很多生涯发展方面的启发。
    
    \item 志同道合的同学。如果你急于解决5.3.2小节中提到的问题（4）“社恐，想社交但实在不会、不敢，怎么办？”，亦或者你在认为自己的满腔热血和抱负需要寻找伙伴一起实现，那么学生会各部门一定能为你提供这么志同道合的一批人，助你实现自己的理想。
\end{itemize}

\vspace{10pt}

学生会现有办公室、外联部、内联部、文体部、宣传部、学术实践部六个职能部门，预计在新生军训、新生训练营期间进行招新，招新对象为工学院全体大一、大二同学。详细介绍请参考公众号“工映青春”上的招新推送。去年的推送为2024.8.21《工学·招募｜工学院学生会招新开始啦！》，敬请期待今年的招新推送。

\vspace{10pt}

随着新工科时代的到来，工学院以及工院学生会的影响力正越来越大，诚挚欢迎大家加入工学院学生会，共创我们的工学新天地！

\vspace{10pt}


\subsection{党建工作委员会}

工学院党建工作委员会为中共北京大学工学院委员会下属工作机构，直接负责基层党组织各类事务，引领党员坚定理想信念，吸纳更多优秀人才加入党组织。二十大以来，构建新时代“系统党建”，充分发挥党员先锋模范作用、基层党组织战斗堡垒作用、人民群众参与作用，是提高党的建设质量的必然选择。党建工作委员会作为党的建设力量的火车头，每一名成员都能够在这里获得思想的淬炼、政治的历练和能力的锻炼。党建工作委员会下设组织部、学习部和综合办公室。

\vspace{10pt}


\subsection{班级建设工作委员会}

工学院班级建设工作委员会是工学院下属工作机构，负责班级各类事务及班主任、辅导员、班委队伍组织建设，旨在夯实工学院基层班级组织建设，引导班级在同学中发挥更大影响力与号召力，加强班主任、带班辅导员、及青年委员等班委队伍凝聚力，提升班级文化建设活动质量。

\section{社团简介}
每学期开学的第二或第三周周末，在百讲与新太阳区域会举行“百团大战”。届时会有上百个社团在场招新，并且在摊位上举办各自的特色小活动。如有兴趣，可以积极参加百团大战并添加相应社团公众号/微信群进行了解。

\vspace{10pt}

社团的一些特点是：

1. 自治性很强，其活动在团委指导下由学生自行组织。

2. 对外开放性不大，导致很多社团甚至不为人所知，看似自娱自乐，实则社员中卧虎藏龙，社团内到底在研发什么“秘密武器”，外人啥也不知道。

3. 加入社团没有任何硬性条件考核，完全出于兴趣，很好地体现了北大“思想自由，兼容并包”的精神。

4. 在社团里好好干还是划水、参加还是不参加活动，完全不受强制，完全出于个人意愿，十分自由。因此一般不会大量占用学习时间导致成绩下滑。

5. 不同社团大小、氛围差异很明显。

6. 社团中也会有负责人等，这也是一种学生工作。

\vspace{10pt}

因此，加不加社团，加几个社团，也是一件“需要你自己衡量”的事。

\section{入党相关}
加入中国共产党是一件极为严肃的事情。同学们应当在坚定信念后再选择入党，而不是人云亦云，“跟风”入党。

\vspace{10pt}

入党有以下几个流程：提交申请书——入党积极分子——发展对象——预备党员——正式党员。以上每次身份的转变都需要经过学习、行动与考察，且周期比较长（总计至少两年多）。

\vspace{10pt}

由于流程很长且较为繁琐，这里不赘述。如果读者有志入党，随时可以提交入党申请书给党支部。注意，入党申请书不必等人统一收取，这是很多人的误区；一般来说给班级辅导员也可以，由其转交。提交申请书后，后续所有流程都会接到通知。

\vspace{10pt}

另外提醒：要牢牢记住重要时间点，如何时入团/党，何日提交申请书，何日被推荐为积极分子等等，这会为以后填写资料节省很多时间。

\section{其他学生组织与课外活动}
\subsection{志愿服务}

志愿服务是提升人奉献精神和工作能力的良好契机。同时，一定的志愿时长\footnote{在“志愿北京”小程序或网站可查看具体时长}还能给综测带来加分，从而给评优评先带来一定优势。因此，花费适当课余时间在志愿服务上，无论从功利还是非功利的角度，都是非常值得的。当然，我们也不鼓励大家去卷过多的志愿时长\footnote{注意综测加分48小时封顶}，重心仍然应该放在课业学习中。

\vspace{10pt}

首先，需要在志愿北京官网 https://www.bv2008.cn/ 上注册账号，并自行记住用户名、密码、志愿者编号等信息。以后在你参加志愿活动后，组织方一般会要求提供志愿者编号，并为你录入相应时长或下发纸质证）。

\vspace{10pt}

有的同学苦恼于志愿时长无处可觅。因此，下面介绍几个重要的志愿服务渠道。具体的信息收集仍需要同学们自力更生。
\begin{itemize}
	\item 校级活动：五四青春长跑志愿者、北京大学图书馆志愿者、返乡宣讲、学生会/团委/青协等学生组织的活动。
	\item 院级活动：较为零散，但是有活动必有志愿服务需求。有意向的同学请咨询学工老师或学长学姐以进入工学院志愿者招募群。
\end{itemize}

\subsection{学校活动}
学校活动数不胜数，种类繁多，只要你想得到，一般就能找得到。

\vspace{10pt}

学校层面，除了基本的学工、社团之外，还有百周年纪念讲堂的演出\footnote{详见百周年纪念讲堂公众号}、佳节时分\footnote{中秋、跨年等}的歌会、十佳歌手大赛、十佳主持人大赛、各种球类的新生杯/北大杯、餐饮中心的厨艺大赛、图书馆的讲座、地方支教等等各种形式的活动，且每个都能带给你别样的体验。活动报名信息获取渠道极为广泛，有意愿的同学请务必时刻关注相关公众号，善用搜索，善于向身边人提问。

\vspace{10pt}

学院层面，值得参与的活动也很多。例如由学生会承办的迎新晚会、保研出国经验分享会等；院友之声系列、学生心理类系列等讲座，在此列举不尽。主要关注的信息渠道有：公众号（见手册7.2.1节）、年级群、班级群等。

\vspace{10pt}

眼观六路，耳听八方，积极获取信息，相信你一定能找到自己热爱的舞台！

\subsection{一则鼓励}
如果你有好的想法，可以向上级建议或自行创办有积极意义的学生组织！
